\begin{table}[htbp]
  \centering
  \caption{Top-bin sensitivity and tests of the quadratic restriction}
  \label{tab:quadratic-preferred-sensitivity}
  \small
  \begin{tabular}{rrrrrr}
    \toprule
    Size for 101+ & Formal (\%) & Informal (\%) & $p_F$ & $p_I$ & $p_{Joint}$ \\[0.25em]
    \midrule
    \multicolumn{6}{l}{\textit{Sector-department clustering}} \\
    101 & 17.9 & 43.2 & $<0.001$ & $<0.001$ & $<0.001$ \\[0.25em]
    150 & 18.9 & 41.9 & $<0.001$ & $<0.001$ & $<0.001$ \\[0.25em]
    250 & 19.8 & 40.5 & $<0.001$ & $<0.001$ & $<0.001$ \\[0.25em]
    \midrule
    \multicolumn{6}{l}{\textit{Sector clustering}} \\
    101 & 17.9 & 43.2 & $<0.001$ & $<0.001$ & $<0.001$ \\[0.25em]
    150 & 18.9 & 41.9 & $<0.001$ & $<0.001$ & $<0.001$ \\[0.25em]
    250 & 19.8 & 40.5 & $<0.001$ & 0.001 & $<0.001$ \\[0.25em]
    \bottomrule
  \end{tabular}
  \vspace{0.3em}
  \begin{minipage}{0.95\textwidth}
  \footnotesize
  Notes: Formal and informal percentages are the implied 101+/solo premiums from the quadratic model, re-estimated for each assigned top-bin size. All other size assignments, observations, controls, and weights remain fixed. The $p$-values test whether the unrestricted categorical profiles from Equation~(\ref{eq:main}) lie on the specified quadratic. There are six restrictions per formality group and twelve jointly. Cluster-robust Wald statistics divided by the number of restrictions use an $F_{q,G-1}$ reference distribution. There are 559 sector-department clusters or 17 sector clusters. These are approximate tests, especially with only 17 clusters. The repeated percentages emphasize that changing clustering does not change point estimates. The value 150 is the baseline assignment, not an observed mean for 101+.
  \end{minipage}
\end{table}
